\chapter{Directives and Symbols}
\label{ch:directives-and-symbols}

\section{Placing Code and Data}
\label{sec:org}

\texttt{.ORG addr} sets the word address at which the following instructions and data are assembled. The address is a
bare decimal or \texttt{0x}-prefixed hexadecimal number, without the hash used for immediate operands. It must be
even, because instructions and data directives each occupy two words. \texttt{.ORG} may appear any number of times, to
put code or data at fixed addresses such as the reset and exception vector table.

% qa:assemble
\begin{lstlisting}
    .ORG 0x0000
    .DQ @start               ; reset vector

    .ORG 0x0200
:start
    NOOP
\end{lstlisting}

The vector table and the address of each vector are defined in the Exception and Interrupt Handling chapter of the
\archmanual{}. Addresses are word addresses, not byte addresses.

\section{Data Directives}

\texttt{.DB}, \texttt{.DW} and \texttt{.DQ} each reserve one slot of two words (4 bytes), the same size as an
instruction, and fill it with a value. Despite their names they differ only in where the value is placed within the
slot:

\begin{table}[H]
	\centering
	\begin{tabular}{lll}
		\toprule
		\textbf{Directive}        & \textbf{Value goes in} & \textbf{Example result}    \\
		\midrule
		\texttt{.DB \#0xAB}       & Low byte               & Bytes \texttt{00 00 00 AB} \\
		\texttt{.DW \#0xABCD}     & Low word               & Bytes \texttt{00 00 AB CD} \\
		\texttt{.DQ \#0xABCD1234} & Whole slot             & Bytes \texttt{AB CD 12 34} \\
		\bottomrule
	\end{tabular}
	\caption{Data directives}
	\label{tab:data-directives}
\end{table}

\texttt{.DB} does not pack four bytes into a slot, and \texttt{.DW} does not pack two words. Choose the directive that
best documents the intended width of the value. Each may also take a label reference or a named constant instead of a
number.

\section{Named Constants}

\texttt{.EQU \$NAME \#value} defines a constant. Write \texttt{\$NAME}, without a hash, wherever the number would be
used. A constant may be used in an immediate operand and in a data directive.

% qa:assemble
\begin{lstlisting}
    .EQU $UART_RECV #0x0001

    MOVI al, $UART_RECV
\end{lstlisting}

\section{Labels and Symbol References}

A label is declared with a leading colon and referred to with a leading at sign. The assembler records the references;
the linker resolves them once it knows where every label ended up, so a reference may be to a label that appears later
in the file or in a different object file.

What a reference resolves to depends on where it is written:

\begin{table}[H]
	\centering
	\begin{tabular}{lp{7cm}}
		\toprule
		\textbf{Where}                                                                                                                                     & \textbf{Default resolution}                                       \\
		\midrule
		\mnemonic{BRAN}, \mnemonic{BRSR}                                                                                                                   & 16-bit signed word displacement from the instruction to the label \\
		Other instruction operands (\mnemonic{MOVI}, \mnemonic{LOAD}, \mnemonic{STOR}, \mnemonic{LDEA}, \mnemonic{LDEL}, \mnemonic{LJMP}, \mnemonic{LJSR}) & The label's low word only                                         \\
		\texttt{.DB}, \texttt{.DW}, \texttt{.DQ}                                                                                                           & The label's full 24-bit address                                   \\
		\bottomrule
	\end{tabular}
	\caption{Default resolution of a label reference}
	\label{tab:symbol-resolution}
\end{table}

A suffix overrides the default:

\begin{table}[H]
	\centering
	\begin{tabular}{ll}
		\toprule
		\textbf{Written as} & \textbf{Resolves to}                   \\
		\midrule
		\texttt{@label.r}   & The PC-relative word displacement      \\
		\texttt{@label.l}   & The low word of the address            \\
		\texttt{@label.u}   & The high word (segment) of the address \\
		\bottomrule
	\end{tabular}
	\caption{Symbol reference suffixes}
	\label{tab:symbol-suffixes}
\end{table}

To load a full 24-bit address into an address register pair, load each half separately:

% qa:assemble
\begin{lstlisting}
    MOVI ah, @table.u         ; segment (high word)
    MOVI al, @table.l         ; offset (low word)
\end{lstlisting}

A displacement that does not fit in 16 signed bits is a link-time error, and so is a reference to a label that is not
declared in any linked object file (Chapter~\ref{ch:diagnostics}).
